\section{随机变量与期望}

前两章谈的都是事件：某个事件发生与否，概率是多少。可多数时候我们关心的不是「发生了什么」，而是「因此得到什么数」——掷两颗骰子，要看的是和，不是两个面；一季雨后要看水库水位，不是每天下了几毫米。把结果换算成一个数，概率论才从「集合上的测度」变成「能加减乘除的东西」；而这一步带来的第一个新对象，就是随机变量。紧接着的两个问题也自然冒出来：这个数平均是多少（期望），它离平均多远（方差、协方差）。本章的顺序就是照这三件事排的。

\subsection{随机变量与分布函数}

\begin{definition}[随机变量]
设样本空间为 $S$。若对每个结果 $\omega\in S$，都有唯一确定的实数 $X(\omega)$ 与之对应，则称 $X=X(\omega)$ 为随机变量。
\end{definition}

既然它的值由结果决定，就可以给它的各个取值赋概率。注意这里换了视角：随机变量本身是一张「把结果换算成数」的表，是确定的函数；随机的是它取到的值。初学时把这两件事混在一起，后面「$X$ 的分布」「$X$ 的期望」就都读不顺了。

\begin{example}[两枚骰子的和]
设 $X$ 为两颗均匀骰子点数之和，则 $X$ 取整数 $2,3,\dots,12$，且
\[
P\{X=2\}=\frac{1}{36},\quad P\{X=3\}=\frac{2}{36},\quad\dots,\quad P\{X=7\}=\frac{6}{36},\quad\dots,\quad P\{X=12\}=\frac{1}{36}.
\]
\end{example}

$X$ 总要取某个值，所以 $1=P(S)=\sum_{i=2}^{12}P\{X=i\}$，这也正是上表能对得上的检验。

有一类随机变量简单到常被忽略，却极其好用：

\begin{definition}[指示变量]
设 $A$ 是事件，定义
\[
I_A(\omega)=\begin{cases}1, & \omega\in A,\\ 0, & \omega\notin A .\end{cases}
\]
称 $I_A$ 为事件 $A$ 的指示随机变量。
\end{definition}

\begin{example}[指示变量的用法]
买两个电子元件，每个或次品或合格，四种结果的概率分别为
\[
(d,d):.09,\quad (d,a):.21,\quad (a,d):.21,\quad (a,a):.49 .
\]
令 $X$ 为合格元件个数，则 $P\{X=0\}=.09$，$P\{X=1\}=.42$，$P\{X=2\}=.49$。若只关心「至少有一个合格」，取
\[
I=\begin{cases}1, & X\ge 1,\\ 0, & X=0,\end{cases}
\]
于是 $P\{I=1\}=.91$，$P\{I=0\}=.09$。
\end{example}

把「是/否」编码成 $1/0$，等于把一个事件换成一个数，后面算期望、算和的方差都会省大力气。

按取值特点分两类。取值能排成一个有限或无穷序列的，叫离散型；在一个区间上连续取值的，叫连续型。取非负整数值的必然是离散的；某辆车在 $(a,b)$ 内任一时点出故障的寿命，是连续的。

\begin{definition}[分布函数]
随机变量 $X$ 的分布函数定义为
\[
F(x)=P\{X\le x\},\qquad x\in\mathbb{R},
\]
记作 $X\sim F$。
\end{definition}

分布函数是「一把抓」的工具：凡关于 $X$ 的概率问题，都能通过它回答。理由只有一步——$\{X\le b\}$ 是 $\{X\le a\}$ 与 $\{a<X\le b\}$ 的不交并，由可加性
\[
P\{X\le b\}=P\{X\le a\}+P\{a<X\le b\}
\quad\Longrightarrow\quad
P\{a<X\le b\}=F(b)-F(a).
\]

\begin{example}
设
\[
F(x)=\begin{cases}0, & x\le 0,\\ 1-e^{-x^2}, & x>0 .\end{cases}
\]
求 $P\{X>1\}$。
\end{example}

用对立事件即可：$P\{X>1\}=1-P\{X\le 1\}=1-F(1)=e^{-1}\approx .368$。

\subsection{分布列与密度}

\begin{definition}[分布列]
设 $X$ 离散，取值为 $x_1,x_2,\dots$，称
\[
p(x_i)=P\{X=x_i\},\qquad i\ge 1
\]
为 $X$ 的分布列（概率质量函数），它满足 $p(x_i)>0$、$p(x)=0$（其余 $x$）、$\sum_i p(x_i)=1$。
\end{definition}

\begin{example}
设 $X$ 取值 $1,2,3$，已知 $p(1)=\frac12$，$p(2)=\frac13$，由概率之和为 $1$ 得 $p(3)=\frac16$。
分布函数可由分布列回推：$F(a)=\sum_{x\le a}p(x)$。
\end{example}

若 $x_1<x_2<x_3<\dots$，则 $F$ 是阶梯函数：在每个区间 $[x_{i-1},x_i)$ 上取常值，在 $x_i$ 处跳起 $p(x_i)$。以 $p(1)=\frac12,\ p(2)=\frac13,\ p(3)=\frac16$ 为例，
\[
F(a)=\begin{cases}
0, & a<1,\\
\frac12, & 1\le a<2,\\
\frac56, & 2\le a<3,\\
1, & a\ge 3 .
\end{cases}
\]
读阶梯函数的诀窍是「只看跳点」：每处跳的高度就是那一点的概率，两个跳点之间不变，整体最终升到 $1$。

连续型则把求和换成积分：

\begin{definition}[概率密度]
设 $X$ 的分布函数 $F$ 可微。若非负函数 $f$ 满足对任意实数集 $B$
\[
P\{X\in B\}=\int_B f(x)\,\mathrm{d}x,
\]
则称 $f$ 为 $X$ 的概率密度函数。
\end{definition}

由「$X$ 必取某值」得 $\int_{-\infty}^{+\infty}f(x)\,\mathrm{d}x=1$；取 $B=[a,b]$ 得 $P\{a\le X\le b\}=\int_a^b f(x)\,\mathrm{d}x$。再令 $a=b$，立刻得到一条初学者常觉得反直觉却极重要的结论：
\[
P\{X=a\}=\int_a^a f(x)\,\mathrm{d}x=0 .
\]
连续型随机变量取任何一个指定值的概率都是零。也正因如此，连续情形下「$<$」和「$\le$」写起来没差别。

由 $F(a)=\int_{-\infty}^{a}f$ 求导得 $\dfrac{\mathrm{d}}{\mathrm{d}a}F(a)=f(a)$：密度是分布函数的导数。至于直观意义，取一个小 $\varepsilon$：
\[
P\Big\{a-\frac{\varepsilon}{2}\le X\le a+\frac{\varepsilon}{2}\Big\}
=\int_{a-\varepsilon/2}^{a+\varepsilon/2}f(x)\,\mathrm{d}x\approx \varepsilon f(a),
\]
也就是
\[
f(a)\approx\frac{\text{$X$ 落在 $a$ 附近的概率}}{\text{区间宽度}} .
\]
所以 $f(a)$ 衡量的是「$X$ 有多容易落在 $a$ 附近」，它本身不是概率，可以大于 $1$——这是第二个坑。

\begin{example}
设 $X$ 的密度为
\[
f(x)=\begin{cases}C(4x-2x^2), & 0<x<2,\\ 0, & \text{其他}.\end{cases}
\]
求 $C$，并计算 $P\{X>1\}$。
\end{example}

由归一化条件，
\[
C\int_0^2(4x-2x^2)\,\mathrm{d}x
=C\Big[2x^2-\frac{2x^3}{3}\Big]_0^2
=C\cdot\frac{8}{3}=1
\ \Longrightarrow\ C=\frac38 .
\]
于是
\[
P\{X>1\}=\frac38\int_1^2(4x-2x^2)\,\mathrm{d}x=\frac12 .
\]
这里没有「$C$ 解出来就完事」这回事：常数靠归一化定，概率靠再积一次，两件事别混成一步。

\subsection{期望}

\begin{definition}[离散型的期望]
设 $X$ 离散，取值为 $x_1,x_2,\dots$，则
\[
E[X]=\sum_i x_i\,P\{X=x_i\}.
\]
\end{definition}

它是可能取值按概率加权的平均。从频率角度看得最清楚：把 $X$ 当作一局游戏的奖金，玩 $N$ 局，取值 $x_i$ 大约出现 $Np(x_i)$ 次，总奖金约 $\sum_i x_iNp(x_i)$，平均每局
\[
\frac{\sum_i x_iNp(x_i)}{N}=\sum_i x_ip(x_i)=E[X].
\]

\begin{example}[掷骰子]
设 $X$ 为均匀骰子的点数，则
\[
E[X]=1\cdot\frac16+2\cdot\frac16+\cdots+6\cdot\frac16=\frac72 .
\]
\end{example}

要提醒一句：$\frac72$ 根本不是 $X$ 能取到的值。期望不是「预期会看到的那个数」，而是大量重复下的平均值。把它当成「最可能的取值」是最常见的误读。

\begin{example}[指示变量]
若 $I$ 是事件 $A$ 的指示变量，则
\[
E[I]=1\cdot P(A)+0\cdot P(A^c)=P(A).
\]
\end{example}

指示变量的期望就是它所指事件的概率。等号右边是概率，左边是数，这个换算后面反复用到。

\begin{example}[多选题的期望得分]
一道单选题四个选项，恰有一个正确。选对得 $2$ 分，不选得 $0$ 分，选错得 $-1$ 分。若完全随机地作答，设 $X$ 为得分，则
\[
E[X]=2\cdot\frac14+(-1)\cdot\frac34=-\frac14 .
\]
\end{example}

期望为负，意思是长期这样瞎选平均是要倒扣分的。期望的符号在这里就是决策依据，这比「期望是个平均数」更有用。

\begin{example}[熵]
一条消息「$X=x$」携带多少信息？用 $I(p)$ 表示概率为 $p$ 的事件被报道时的信息量，它应当非负、且随 $p$ 增大而减小。若 $X,Y$ 独立，$P\{X=x\}=p$，$P\{Y=y\}=q$，两条消息合起来的信息是 $I(p)+I(q)$；而联合事件的概率是 $pq$，故
\[
I(pq)=I(p)+I(q).
\]
令 $G(p)=I(2^{-p})$，上式变成 $G(p+q)=G(p)+G(q)$，单调解只有 $G(p)=cp$。于是 $I(q)=-c\log_2 q$，取 $c=1$，单位是比特。
\end{example}

按这个思路，$X$ 传递的期望信息量就是熵
\[
H(X)=-\sum_{i=1}^{n}p_i\log_2 p_i .
\]
熵没有在本课程里展开，但它是一个纯由期望定义出来的量，正好说明期望不只是「平均分」。

连续情形把求和换成积分：

\begin{definition}[连续型的期望]
设 $X$ 的密度为 $f$，则
\[
E[X]=\int_{-\infty}^{+\infty}xf(x)\,\mathrm{d}x .
\]
\end{definition}

理由仍是加权平均：当 $\mathrm{d}x$ 很小时 $f(x)\,\mathrm{d}x\approx P\{x<X<x+\mathrm{d}x\}$，于是「取值 $\times$ 概率」逐段累加就成了积分。

\begin{example}[等消息]
下午五点后 $X$ 小时收到一条消息，$X$ 的密度为 $f(x)=1/1.5$（$0<x<1.5$）。则
\[
E[X]=\int_0^{1.5}\frac{x}{1.5}\,\mathrm{d}x=.75,
\]
平均要等三刻钟。
\end{example}

把期望想成重心最能帮助记忆：在无重量的杆上，于 $x_i$ 处挂质量 $p(x_i)$，平衡点就是 $E[X]$。验算一下力矩：
\[
\sum_i\big(x_i-E[X]\big)p(x_i)=E[X]-E[X]=0 .
\]
两个附带的性质：期望和 $X$ 同量纲（概率没单位，加权不会改变单位）；期望可以不落在取值集合里，重心落在两枚砝码之间是常事。

\subsection{随机变量函数的期望}

要算 $E[g(X)]$，笨办法是先求 $g(X)$ 的分布再按定义算。有一条短得多的路：

\begin{proposition}
\label{prop:lotus}
\begin{enumerate}
\item 若 $X$ 离散、分布列为 $p$，则对任意实值函数 $g$
\[
E[g(X)]=\sum_x g(x)p(x);
\]
\item 若 $X$ 连续、密度为 $f$，则
\[
E[g(X)]=\int_{-\infty}^{+\infty}g(x)f(x)\,\mathrm{d}x .
\]
\end{enumerate}
\end{proposition}

要点是：加在 $g(x)$ 上的权重仍然是「$X=x$」的概率（或概率密度），不需要先知道 $g(X)$ 长什么样。

\begin{example}
设 $p(0)=.2,\ p(1)=.5,\ p(2)=.3$，求 $E[X^2]$。
\end{example}

绕路算：令 $Y=X^2$，取值 $0,1,4$，概率分别 $.2,.5,.3$，故 $E[Y]=0(.2)+1(.5)+4(.3)=1.7$。
直接用命题~\ref{prop:lotus}：$E[X^2]=0^2(.2)+1^2(.5)+2^2(.3)=1.7$。两条路当然一致，但第二条省掉了求分布这一步。

\begin{example}[维修成本]
某厂查找并修好一次电气故障所需时间 $X$（小时）的密度为
\[
f_X(x)=\begin{cases}1, & 0<x<1,\\ 0, & \text{其他},\end{cases}
\]
停机 $x$ 小时的损失为 $x^3$。求期望损失。
\end{example}

先绕路：$Y=X^3$，对 $0\le a\le 1$ 有 $F_Y(a)=P\{X\le a^{1/3}\}=a^{1/3}$，故 $f_Y(a)=\frac13a^{-2/3}$，
\[
E[X^3]=\int_0^1 a\cdot\frac13a^{-2/3}\,\mathrm{d}a=\frac13\cdot\frac34=\frac14 .
\]
直接算：$E[X^3]=\int_0^1x^3\,\mathrm{d}x=\frac14$。同样一条结论，后者省了求导和换元。

\begin{corollary}
对常数 $a,b$，
\[
E[aX+b]=aE[X]+b .
\]
\end{corollary}

离散情形一行就完：
\[
\sum_x(ax+b)p(x)=a\sum_x xp(x)+b\sum_xp(x)=aE[X]+b,
\]
连续情形把求和换成积分即可。取 $a=0$ 得 $E[b]=b$（常数的期望是它自己）；取 $b=0$ 得 $E[aX]=aE[X]$。

$E[X]$ 又叫均值或一阶矩。一般地，$n\ge 1$ 时称
\[
E[X^n]=\begin{cases}\sum_x x^np(x), & X\ \text{离散},\\[2pt] \displaystyle\int_{-\infty}^{+\infty}x^nf(x)\,\mathrm{d}x, & X\ \text{连续}\end{cases}
\]
为 $n$ 阶矩。

\subsection{和的期望}

命题~\ref{prop:lotus} 的二维版本是
\[
E[g(X,Y)]=\sum_y\sum_x g(x,y)p(x,y)
\qquad\text{或}\qquad
\int_{-\infty}^{+\infty}\!\!\int_{-\infty}^{+\infty}g(x,y)f(x,y)\,\mathrm{d}x\,\mathrm{d}y .
\]
取 $g(x,y)=x$ 与 $g(x,y)=y$ 分别代入再相加，就得到下面这条极有用的性质。

\begin{proposition}[期望的可加性]
对任意随机变量 $X_1,\dots,X_n$——不论是否独立——
\[
E[X_1+X_2+\cdots+X_n]=E[X_1]+E[X_2]+\cdots+E[X_n].
\]
\end{proposition}

没有独立性条件，这条比它看起来更值钱。下面几个例子都能说明它顶多少计算量。

\begin{example}[两颗骰子之和]
两颗骰子之和 $X$ 的期望，不必去算 $\sum_{i=2}^{12}iP\{X=i\}$。写成 $X=X_1+X_2$，则
\[
E[X]=E[X_1]+E[X_2]=\frac72+\frac72=7 .
\]
\end{example}

\begin{example}[三份投标]
某公司投标三个项目，利润分别为 $10,20,40$（千元），中标概率为 $.2,.8,.3$，求期望总利润。
\end{example}

设 $X_i$ 为第 $i$ 个项目的利润，则 $E[X_1]=2$，$E[X_2]=16$，$E[X_3]=12$，总期望 $30$ 千元。注意这里没有假设独立性——三个项目同不同时中，都不影响和的期望。

\begin{example}[匹配问题]
把 $N$ 封信随机塞进 $N$ 个写好地址的信封，每封信进每个信封的机会相同。求装对的信封数的期望。
\end{example}

设 $X_i$ 为「第 $i$ 封信装对」的指示变量，$X=\sum_{i=1}^NX_i$。每封信进对信封的概率是 $1/N$，故 $E[X_i]=1/N$，于是
\[
E[X]=\sum_{i=1}^NE[X_i]=N\cdot\frac1N=1 .
\]
不管多少封信，平均恰好装对一封。这里 $X_i$ 之间当然不独立（一封装对会影响别的），但可加性不要求独立——这正是它好用的地方。

\begin{example}[优惠券种类]
有 $20$ 种优惠券，每张等可能地是其中任一种。随机取 $10$ 张，期望能集到多少种？
\end{example}

设 $X_i$ 为「第 $i$ 种出现过」的指示变量，$X=\sum_{i=1}^{20}X_i$。每张不是第 $i$ 种的概率是 $19/20$，$10$ 张都不是的概率是 $(19/20)^{10}$，故
\[
E[X_i]=P\{X_i=1\}=1-\Big(\frac{19}{20}\Big)^{10},
\qquad
E[X]=20\Big[1-\Big(\frac{19}{20}\Big)^{10}\Big]\approx 8.025 .
\]
十张券平均能凑到八种左右。这类「至少一个」的问题，用指示变量把并集拆成和，是标准解法。

期望还有一个变分刻画，说明了它在什么意义上「最代表」$X$：

\begin{proposition}[均值是最佳常数预测]
设 $\mu=E[X]$。对任意常数 $c$，
\[
E[(X-c)^2]\ge E[(X-\mu)^2],
\]
等号成立当且仅当 $c=\mu$。
\end{proposition}

证明只有配方一步：
\[
E[(X-c)^2]=E[(X-\mu+\mu-c)^2]=E[(X-\mu)^2]+2(\mu-c)E[X-\mu]+(\mu-c)^2,
\]
而 $E[X-\mu]=0$，所以它等于 $E[(X-\mu)^2]+(\mu-c)^2\ge E[(X-\mu)^2]$。

为什么用平方误差而不用 $E[X-c]$？因为原始误差会互相抵消。取 $X$ 以各 $\frac12$ 的概率等于 $+100$ 或 $-100$，预测 $c=0$ 时平均误差
\[
E[X-0]=.5(100)+.5(-100)=0,
\]
看起来完美，可每次预测都差了 $100$。所以评价预测好坏时，必须先把误差变成非负的量。

\subsection{方差}

三个随机变量
\[
W\equiv 0,\qquad
Y=\begin{cases}-1, & \text{概率 }1/2\\ 1, & \text{概率 }1/2\end{cases},\qquad
Z=\begin{cases}-100, & \text{概率 }1/2\\ 100, & \text{概率 }1/2\end{cases}
\]
的期望都是 $0$，可散布程度天差地别。期望只描述位置，不描述形状，所以还需要第二个量。用 $E[|X-\mu|]$ 也说得通，但绝对值处理起来麻烦，平方可导可积，于是取平方。

\begin{definition}[方差与标准差]
设 $X$ 的均值为 $\mu$，称
\[
\mathrm{Var}(X)=E\big[(X-\mu)^2\big]
\]
为 $X$ 的方差；$\sqrt{\mathrm{Var}(X)}$ 称为标准差，与 $X$ 同量纲。
\end{definition}

实际计算几乎总用下面这个变形：
\[
\mathrm{Var}(X)=E[X^2-2\mu X+\mu^2]=E[X^2]-2\mu E[X]+\mu^2=E[X^2]-\big(E[X]\big)^2 .
\]

\begin{example}[掷骰子]
均匀骰子 $E[X^2]=\frac16(1^2+2^2+\cdots+6^2)=\frac{91}{6}$，而 $E[X]=\frac72$，于是
\[
\mathrm{Var}(X)=\frac{91}{6}-\Big(\frac72\Big)^2=\frac{35}{12}.
\]
\end{example}

\begin{example}
对 $p(0)=.2,\ p(1)=.5,\ p(2)=.3$，前面已算出 $E[X]=1.1$、$E[X^2]=1.7$，故
\[
\mathrm{Var}(X)=1.7-(1.1)^2=0.49 .
\]
\end{example}

\begin{example}[指示变量的方差]
设 $I$ 是事件 $A$ 的指示变量。由 $I^2=I$，
\[
\mathrm{Var}(I)=E[I^2]-\big(E[I]\big)^2=E[I]\big(1-E[I]\big)=P(A)\big[1-P(A)\big].
\]
\end{example}

顺带一条不等式：因为 $(X-\mu)^2\ge 0$，故 $\mathrm{Var}(X)\ge 0$，也就是
\[
E[X^2]\ge\big(E[X]\big)^2 .
\]
随机变量的平方期望不小于期望的平方，等号只在 $X$ 几乎为常数时成立。

\begin{example}[友谊悖论]
把「平均来说，你的朋友比你朋友更多」写成一个可证的不等式。设第 $i$ 个人的朋友数为 $f(i)$，总度数 $\bar f=\sum_if(i)$。均匀随机选一个人 $X$，则 $P\{X=i\}=1/n$，故
\[
E[f(X)]=\sum_{i=1}^n\frac{f(i)}{n}=\frac{\bar f}{n}.
\]
换个抽法：让每个人把每位朋友的名字写在一张纸条上，共 $\bar f$ 张，随机抽一张，记上面的名字为 $Y$。此时 $P\{Y=i\}=f(i)/\bar f$，于是
\[
E[f(Y)]=\sum_i\frac{f^2(i)}{\bar f}=\frac{E[f^2(X)]}{E[f(X)]}\ \ge\ E[f(X)] ,
\]
最后一步用的正是上面那条不等式。
\end{example}

原因一句话：$X$ 是均匀抽的，$Y$ 是按朋友数加权抽的，人缘好的人更容易被抽到，于是「朋友的朋友数」被系统性地拉高。

\begin{proposition}[仿射变换的方差]
对常数 $a,b$，
\[
\mathrm{Var}(aX+b)=a^2\,\mathrm{Var}(X).
\]
\end{proposition}

三个特例值得记住：$a=0$ 时 $\mathrm{Var}(b)=0$，常数的方差为零；$a=1$ 时 $\mathrm{Var}(X+b)=\mathrm{Var}(X)$，平移不改变散布；$b=0$ 时 $\mathrm{Var}(aX)=a^2\mathrm{Var}(X)$，注意是平方。用机械学的说法：期望是重心，方差是转动惯量。

\subsection{协方差与和的方差}

期望可加，方差不可加。取 $X+X$ 就能看出：
\[
\mathrm{Var}(2X)=4\,\mathrm{Var}(X)\ne \mathrm{Var}(X)+\mathrm{Var}(X).
\]
独立性下它会恢复可加，但要证明这一点，得先引入一个新的量。

\begin{definition}[协方差]
设 $X,Y$ 的均值分别为 $\mu_x,\mu_y$，称
\[
\mathrm{Cov}(X,Y)=E\big[(X-\mu_x)(Y-\mu_y)\big]
\]
为它们的协方差。
\end{definition}

展开乘积即得
\[
\mathrm{Cov}(X,Y)=E[XY]-\mu_xE[Y]-\mu_yE[X]+\mu_x\mu_y=E[XY]-E[X]E[Y].
\]

由定义直接看出三条：$\mathrm{Cov}(X,Y)=\mathrm{Cov}(Y,X)$；$\mathrm{Cov}(X,X)=\mathrm{Var}(X)$；$\mathrm{Cov}(aX,Y)=a\,\mathrm{Cov}(X,Y)$。再往下是双线性：

\begin{lemma}
\[
\mathrm{Cov}(X_1+X_2,Y)=\mathrm{Cov}(X_1,Y)+\mathrm{Cov}(X_2,Y).
\]
\end{lemma}

证明还是展开：
\[
\mathrm{Cov}(X_1+X_2,Y)=E[(X_1+X_2)Y]-E[X_1+X_2]E[Y]
=\mathrm{Cov}(X_1,Y)+\mathrm{Cov}(X_2,Y).
\]
据此对两个变量分别归纳，得到

\begin{proposition}[双线性]
\[
\mathrm{Cov}\Big(\sum_{i=1}^nX_i,\ \sum_{j=1}^mY_j\Big)=\sum_{i=1}^n\sum_{j=1}^m\mathrm{Cov}(X_i,Y_j).
\]
\end{proposition}

\begin{corollary}[和的方差]
\[
\mathrm{Var}\Big(\sum_{i=1}^nX_i\Big)=\sum_{i=1}^n\mathrm{Var}(X_i)+\sum_{i=1}^n\sum_{j\ne i}\mathrm{Cov}(X_i,X_j).
\]
特别地，$n=2$ 时
\[
\mathrm{Var}(X+Y)=\mathrm{Var}(X)+\mathrm{Var}(Y)+2\,\mathrm{Cov}(X,Y).
\]
\end{corollary}

矩阵式地看，$\mathrm{Var}(\sum_iX_i)=\mathrm{Cov}(\sum_iX_i,\sum_jX_j)=\sum_{i,j}\mathrm{Cov}(X_i,X_j)$，对角线上的项就是各个 $\mathrm{Var}(X_i)$。

\begin{theorem}[独立则协方差为零]
若 $X$ 与 $Y$ 独立，则 $\mathrm{Cov}(X,Y)=0$；从而对两两独立的 $X_1,\dots,X_n$，
\[
\mathrm{Var}\Big(\sum_{i=1}^nX_i\Big)=\sum_{i=1}^n\mathrm{Var}(X_i).
\]
\end{theorem}

只需证 $E[XY]=E[X]E[Y]$。离散情形：
\[
E[XY]=\sum_j\sum_ix_iy_jP\{X=x_i,Y=y_j\}
=\sum_jy_jP\{Y=y_j\}\sum_ix_iP\{X=x_i\}=E[Y]E[X],
\]
中间一步用的正是独立性。连续情形同理。

\begin{remark}
反过来不成立：协方差为零推不出独立。这一点必须记住——定理只能单方向使用。
\end{remark}

\begin{example}[零协方差但不独立]
设 $X$ 在 $\{-1,0,1\}$ 上等概率取值，$Y=X^2$。则
\[
E[X]=\frac{-1+0+1}{3}=0,\qquad
E[Y]=\frac{1+0+1}{3}=\frac23,\qquad
E[XY]=E[X^3]=0,
\]
故
\[
\mathrm{Cov}(X,Y)=0-0\cdot\frac23=0 .
\]
但
\[
P\{X=1,Y=1\}=\frac13\ne P\{X=1\}P\{Y=1\}=\frac13\cdot\frac23=\frac29,
\]
所以不独立。$Y$ 完全由 $X$ 决定，依赖强得很，只是正负两半抵消，线性相关为零。
\end{example}

\begin{example}[十次掷骰子之和]
设 $X_i$ 为第 $i$ 次掷骰子的点数，独立同分布，则
\[
\mathrm{Var}\Big(\sum_{i=1}^{10}X_i\Big)=10\cdot\frac{35}{12}=\frac{175}{6}.
\]
\end{example}

\begin{example}[十次抛硬币的正面数]
设正面数为 $\sum_{j=1}^{10}I_j$，其中 $I_j$ 为第 $j$ 次为正面的指示变量。由前面指示变量的方差公式，
\[
\mathrm{Var}(I_j)=\frac12\Big(1-\frac12\Big)=\frac14,
\]
于是总方差为 $10\cdot\frac14=\frac{10}{4}$。
\end{example}

协方差的符号有明确含义。取 $X,Y$ 分别是事件 $A,B$ 的指示变量，则 $XY$ 指示 $A\cap B$，
\[
\mathrm{Cov}(X,Y)=P\{X=1,Y=1\}-P\{X=1\}P\{Y=1\}
\]
所以
\[
\mathrm{Cov}(X,Y)>0
\iff
\frac{P\{X=1,Y=1\}}{P\{X=1\}}>P\{Y=1\}
\iff
P\{Y=1\mid X=1\}>P\{Y=1\}.
\]
正协方差说明 $X=1$ 让 $Y=1$ 更容易发生，负的则相反。一般的量纲下不那么干净，所以再除以各自的尺度：

\begin{definition}[相关系数]
\[
\mathrm{Corr}(X,Y)=\frac{\mathrm{Cov}(X,Y)}{\sqrt{\mathrm{Var}(X)\mathrm{Var}(Y)}}\in[-1,1],
\]
它是无量的，衡量的是线性相关程度。
\end{definition}

最后一个例子说明「边缘相同，总量风险不同」。设 $X,Y$ 都在 $\{1,2,3,4,5,6\}$ 上均匀，固定两个边缘、只改联合分布：

\begin{example}[耦合]
\begin{center}
\begin{tabular}{lll}
\hline
联合分布 & $\mathrm{Var}(X+Y)$ & $P\{X+Y\ge 10\}$\\
\hline
$Y=7-X$（完全负相关） & $0$ & $0$\\
$X,Y$ 独立 & $35/6$ & $1/6$\\
$Y=X$（完全正相关） & $35/3$ & $1/3$\\
\hline
\end{tabular}
\end{center}
\end{example}

单个风险一模一样，合起来的风险可以从 $0$ 到 $35/3$。缺的那块信息就是协方差：$\mathrm{Var}(X+Y)=\mathrm{Var}(X)+\mathrm{Var}(Y)+2\mathrm{Cov}(X,Y)$，其余都固定时，和的方差完全由协方差决定。这也是「分散投资」的数学理由。

本章容易踩的地方集中在几处。第一，把期望当成「最可能出现的值」或「一定取得到的值」——骰子的 $7/2$ 就是反例。第二，把密度当成概率，忘了它可以超过 $1$、而某个点的概率恒为 $0$。第三，以为期望可加就需要独立（不需要），又以为方差可加也不需要条件（需要协方差项为零）。第四，把「协方差为零」当成「独立」——最后那个 $Y=X^2$ 的例子专门堵这个口子。最后，算 $\mathrm{Var}(aX)$ 时漏掉平方，是手滑里最常见的一种。
